\answerpaper{数学 B 类}

\answerproblem{一、}{15}
设球面 $S: x^2+y^2+z^2=1$，求以点 $M_0(0,0,a)$（$a\in\R$，$|a|>1$）为顶点的与 $S$ 相切的锥面方程．
\begin{solution}
\textbf{解法一：}

设 $L$ 为过顶点 $M_0(0,0,a)$，方向为 $\vec s=(l,m,n)$，与 $S$ 相切的锥面上的任意一条母线，则对于 $L$ 上任意一点 $M(x,y,z)$，$L$ 的方程可以表示为
\[
  \frac{x-0}{l}=\frac{y-0}{m}=\frac{z-a}{n},
\]
其中 $l,m,n\in\R$ 不全为零．设 $L$ 的参数方程为
\begin{align*}
  \begin{cases}
    x=lt,\\
    y=mt,\\
    z=a+nt,
  \end{cases}\tag{1}
\end{align*}
其中 $t\in\R$ 为参数．

将直线的参数方程 (1) 代入 $S$ 中可得
\[
  (l^2+m^2+n^2)t^2+2ant+a^2-1=0.
\]
由直线 $L$ 与球面 $S$ 相切的条件可知
\[
  (2an)^2-4(l^2+m^2+n^2)(a^2-1)=0,
\]
亦即
\begin{align*}
  (l^2+m^2)a^2=(l^2+m^2+n^2).\tag{2}
\end{align*}
由 (1) 和 (2) 消去参数 $t$ 可得锥面方程
\[
  (a^2-1)(x^2+y^2)-(z-a)^2=0,
\]
其中 $|a|>1$．

\textbf{解法二：}

设 $O(0,0,0)$ 为球心坐标，$M(x,y,z)$ 为切锥面与球面的切点，半顶角为
\[
  \alpha=\angle\bigl(\overrightarrow{M_0O},\overrightarrow{M_0M}\bigr),
\]
则有 $\sin\alpha=\dfrac{1}{|a|}$．

注意到
\[
  \cos^2\alpha=\frac{\abs{\overrightarrow{M_0O}\cdot\overrightarrow{M_0M}}^2}{\abs{\overrightarrow{M_0O}}^2\abs{\overrightarrow{M_0M}}^2},\qquad
  \cos^2\alpha=1-\sin^2\alpha,
\]
得到
\[
  \frac{(z-a)^2}{x^2+y^2+(z-a)^2}=\frac{a^2-1}{a^2},
\]
即
\[
  (a^2-1)(x^2+y^2)-(z-a)^2=0,
\]
其中 $|a|>1$．
\end{solution}

\answerproblem{二、}{15}
设 $B\subset\R^n\ (n\ge2)$ 是单位开球，函数 $u,v$ 在 $\bar B$ 上连续，在 $B$ 内二阶连续可导，满足
\[
\begin{gathered}
  -\Delta u-(1-u^2-v^2)u=0,\qquad x\in B,\\
  -\Delta v-(1-u^2-v^2)v=0,\qquad x\in B,\\
  u(x)=v(x)=0,\qquad x\in\partial B,
\end{gathered}
\]
其中 $x=(x_1,x_2,\cdots,x_n)$，$\Delta u=\dfrac{\partial^2u}{\partial x_1^2}+\dfrac{\partial^2u}{\partial x_2^2}+\cdots+\dfrac{\partial^2u}{\partial x_n^2}$，$\partial B$ 表示 $B$ 的边界．证明：$u^2(x)+v^2(x)\le1\ (x\in\bar B)$．
\begin{solution}
记 $w=w(x)=u^2(x)+v^2(x)$，则 $w$ 满足问题
\begin{align*}
  -\Delta w-2(1-w)w&=-2\bigl(|\nabla u|^2+|\nabla v|^2\bigr),\qquad x\in B,\tag{1}\\
  w(x)&=0,\qquad x\in\partial B.
\end{align*}

显然，$w(x)\in C^2(B)\cap C(\bar B)$．所以，$w(x)$ 必然在 $\bar B$ 上达到最大值．设最大值点为 $x_1$．

若 $x_1\in B$，则 $\nabla w(x_1)=0$，$-\Delta w(x_1)\ge0$．于是由 (1) 得到，在 $x_1$ 处，
\[
  0\le-\Delta w\le2(1-w)w-2\bigl(|\nabla u|^2+|\nabla v|^2\bigr)\le2(1-w)w.
\]
而 $w(x_1)\ge0$，故上式表明 $w(x_1)\le1$．

若 $x_1\in\partial B$，则由 (1)，$w(x_1)=0$．

综合之，恒有 $0\le w\le1$，$x\in\bar B$．
\end{solution}

\answerproblem{三、}{15}
设 $f(x)=x^{2021}+a_{2020}x^{2020}+a_{2019}x^{2019}+\cdots+a_2x^2+a_1x+a_0$ 为整系数多项式，$a_0\ne0$．设对任意 $0\le k\le2020$ 有 $|a_k|\le40$，证明：$f(x)=0$ 的根不可能全为实数．
\begin{solution}
设 $f(x)=0$ 的 2021 个根分别为 $x_1,x_2,\cdots,x_{2021}$．由于 $a_0\ne0$，所以 $x_i\ne0$，$1\le i\le2021$．若 $x_1,x_2,\cdots,x_{2021}$ 都是实数，由 Cauchy 不等式有
\begin{align*}
  \Bigl(\sum_{i=1}^{2021}x_i^2\Bigr)\Bigl(\sum_{i=1}^{2021}\frac{1}{x_i^2}\Bigr)\ge\Bigl(\sum_{i=1}^{2021}x_i\cdot\frac{1}{x_i}\Bigr)^2=2021^2.
\end{align*}
由 Vieta 定理，
\begin{align*}
  \sum_{i=1}^{2021}x_i=-a_{2020},\qquad
  \sum_{1\le i<j\le2021}x_ix_j=a_{2019},
\end{align*}
由此得到
\begin{align*}
  \sum_{i=1}^{2021}x_i^2=\Bigl(\sum_{i=1}^{2021}x_i\Bigr)^2-2\sum_{1\le i<j\le2021}x_ix_j=a_{2020}^2-2a_{2019}.
\end{align*}
注意到 $\dfrac{1}{x_1},\dfrac{1}{x_2},\cdots,\dfrac{1}{x_{2021}}$ 是多项式
\[
  g(x)=x^{2021}f\Bigl(\frac{1}{x}\Bigr)=a_0x^{2021}+a_1x^{2020}+a_2x^{2019}+\cdots+a_{2019}x^2+a_{2020}x+1
\]
的根．继续由 Vieta 定理，
\begin{align*}
  \sum_{i=1}^{2021}\frac{1}{x_i}=-\frac{a_1}{a_0},\qquad
  \sum_{1\le i<j\le2021}\frac{1}{x_i}\cdot\frac{1}{x_j}=\frac{a_2}{a_0},
\end{align*}
所以
\begin{align*}
  \sum_{i=1}^{2021}\frac{1}{x_i^2}=\Bigl(\sum_{i=1}^{2021}\frac{1}{x_i}\Bigr)^2-2\sum_{1\le i<j\le2021}\frac{1}{x_i}\cdot\frac{1}{x_j}=\frac{a_1^2}{a_0^2}-\frac{2a_2}{a_0}.
\end{align*}
因为对任意 $0\le k\le2020$ 有 $|a_k|\le40$，又 $a_0$ 为非零整数，故 $|a_0|\ge1$，所以
\begin{align*}
  \Bigl(\sum_{i=1}^{2021}x_i^2\Bigr)\Bigl(\sum_{i=1}^{2021}\frac{1}{x_i^2}\Bigr)&=\bigl(a_{2020}^2-2a_{2019}\bigr)\Bigl(\frac{a_1^2}{a_0^2}-\frac{2a_2}{a_0}\Bigr)\\
  &\le(40^2+2\cdot40)(40^2+2\cdot40)=1680^2,
\end{align*}
矛盾．证毕．
\end{solution}

\answerproblem{四、}{20}
设 $R=\{0,1,-1\}$，$S$ 为 $R$ 上的 3 阶行列式全体，即 $S=\{\det(a_{ij})_{3\times3}\mid a_{ij}\in R\}$．证明：$S=\{-4,-3,-2,-1,0,1,2,3,4\}$．
\begin{solution}
首先，通过直接检验可知
\[
\begin{gathered}
  \begin{vmatrix}0&0&0\\0&0&0\\0&0&0\end{vmatrix}=0,\qquad
  \begin{vmatrix}1&0&0\\0&1&0\\0&0&1\end{vmatrix}=1,\qquad
  \begin{vmatrix}1&0&0\\0&1&-1\\0&1&1\end{vmatrix}=2,\\
  \begin{vmatrix}1&-1&0\\1&1&-1\\0&1&1\end{vmatrix}=3,\qquad
  \begin{vmatrix}1&1&1\\1&-1&1\\1&1&-1\end{vmatrix}=4.
\end{gathered}
\]

其次，由于交换两行，行列式值改变符号，因此有 $S\supseteq\{-4,-3,-2,-1,0,1,2,3,4\}$．

第三，我们证明：$\forall(a_{ij})_{3\times3}$，$a_{ij}\in R$，总有 $\abs{\det(a_{ij})}\le4$．

事实上，由对角线法则可知
\begin{align*}
  \det(a_{ij})={}&a_{11}a_{22}a_{33}+a_{12}a_{23}a_{31}+a_{13}a_{32}a_{21}\\
  &-a_{13}a_{22}a_{31}-a_{12}a_{21}a_{33}-a_{11}a_{32}a_{23}.
\end{align*}
记
\begin{align*}
  b_1=a_{11}a_{22}a_{33},\quad b_2=a_{12}a_{23}a_{31},\quad b_3=a_{13}a_{32}a_{21},\\
  b_4=-a_{13}a_{22}a_{31},\quad b_5=-a_{12}a_{21}a_{33},\quad b_6=-a_{11}a_{32}a_{23}.
\end{align*}
直接观察可知：每个 $a_{ij}$ 在单项 $b_1,b_2,b_3,b_4,b_5,b_6$ 中共出现两次，且
\begin{align*}
  b_1b_2b_3b_4b_5b_6=-a_{11}^2a_{12}^2a_{13}^2a_{21}^2a_{22}^2a_{23}^2a_{31}^2a_{32}^2a_{33}^2.
\end{align*}
因此立即可得：若有某个 $a_{ij}=0$，则 $b_1,\cdots,b_6$ 中至少有两个为 $0$，从而 $\abs{\det(a_{ij})}\le4$．倘若每个 $a_{ij}$ 都不等于 $0$，则由 $a_{ij}=\pm1$ 得 $b_1,\cdots,b_6$ 之积 $=-1$，从而至少有一个 $b_i$ 为 $-1$，同时也至少有一个 $b_j$ 为 $1$，否则与 $b_1b_2b_3b_4b_5b_6=-1$ 矛盾．结果 $b_i$ 与 $b_j$ 互相抵消，仍有 $\abs{\det(a_{ij})}\le4$．

至此，综上所得，$S=\{-4,-3,-2,-1,0,1,2,3,4\}$，$S$ 共由 $9$ 个元素所组成．
\end{solution}

\answerproblem{五、}{15}
设函数 $f$ 在 $[-1,1]$ 内有定义，在 $x=0$ 的某邻域内连续可导，且 $\lim\limits_{x\to0}\dfrac{f(x)}{x}=a>0$．证明：级数 $\sum\limits_{n=1}^{\infty}(-1)^nf\Bigl(\dfrac{1}{n}\Bigr)$ 收敛，$\sum\limits_{n=1}^{\infty}f\Bigl(\dfrac{1}{n}\Bigr)$ 发散．
\begin{solution}
由 $\lim\limits_{x\to0^+}\dfrac{f(x)}{x}=a>0$，知 $\lim\limits_{x\to0^+}f(x)=0$．又 $f(x)$ 在 $x=0$ 的某邻域内连续可导，则
\[
  f(0)=\lim_{x\to0^+}f(x)=0.
\]
于是
\[
  0<a=\lim_{x\to0^+}\frac{f(x)}{x}=\lim_{x\to0^+}\frac{f(x)-f(0)}{x}=f'(0).
\]
由于 $f'(x)$ 在 $x=0$ 的某邻域内连续，存在正数 $\delta>0$，使得 $\forall x\in[0,\delta]$，有 $f'(x)>0$．因此，在 $[0,\delta]$ 上 $f(x)$ 单调增加．

于是存在正整数 $N>\dfrac{1}{\delta}$，当 $n>N$ 时，$f\Bigl(\dfrac{1}{n}\Bigr)>0$ 且 $f\Bigl(\dfrac{1}{n}\Bigr)>f\Bigl(\dfrac{1}{n+1}\Bigr)$．由 $\lim\limits_{x\to0^+}f(x)=0$，知 $\lim\limits_{n\to\infty}f\Bigl(\dfrac{1}{n}\Bigr)=0$，且 $\sum\limits_{n=N}^{\infty}(-1)^nf\Bigl(\dfrac{1}{n}\Bigr)$ 为交错级数，由莱布尼兹判别法，级数 $\sum\limits_{n=N}^{\infty}(-1)^nf\Bigl(\dfrac{1}{n}\Bigr)$ 收敛，故级数 $\sum\limits_{n=1}^{\infty}(-1)^nf\Bigl(\dfrac{1}{n}\Bigr)$ 收敛．

由于级数 $\sum\limits_{n=1}^{\infty}\dfrac{1}{n}$ 发散，而 $\lim\limits_{n\to\infty}\dfrac{f\bigl(\frac{1}{n}\bigr)}{1/n}=a>0$，故 $\sum\limits_{n=1}^{\infty}f\Bigl(\dfrac{1}{n}\Bigr)$ 发散．
\end{solution}

\answerproblem{六、}{20}
设 $f(x)=\ln\sum\limits_{n=1}^{\infty}\dfrac{\ee^{nx}}{n^2}$．证明函数 $f$ 在 $(-\infty,0)$ 内为严格凸的，并且对任意 $\xi\in(-\infty,0)$，存在 $x_1,x_2\in(-\infty,0)$ 使得
\[
  f'(\xi)=\frac{f(x_2)-f(x_1)}{x_2-x_1}.
\]
（称 $(a,b)$ 内的函数 $f$ 为严格凸的，如果对任何 $\alpha\in(0,1)$，以及 $x,y\in(a,b)$，$x\ne y$ 成立 $f(\alpha x+(1-\alpha)y)<\alpha f(x)+(1-\alpha)f(y)$．）
\begin{solution}
记 $g(x)=\sum\limits_{n=1}^{\infty}\dfrac{\ee^{nx}}{n^2}$．我们有
\[
  f'(x)=\frac{g'(x)}{g(x)},\qquad
  f''(x)=\frac{g''(x)g(x)-(g'(x))^2}{g^2(x)}.
\]
又因为
\[
  g'(x)=\sum_{n=1}^{\infty}\frac{\ee^{nx}}{n},\qquad
  g''(x)=\sum_{n=1}^{\infty}\ee^{nx}.
\]
于是，由 Hölder 不等式，
\begin{align*}
  g''(x)g(x)-(g'(x))^2&=\Bigl(\sum_{n=1}^{\infty}\ee^{nx}\Bigr)\Bigl(\sum_{n=1}^{\infty}\frac{\ee^{nx}}{n^2}\Bigr)-\Bigl(\sum_{n=1}^{\infty}\frac{\ee^{nx}}{n}\Bigr)^2>0,
\end{align*}
从而函数 $f(x)$ 为严格凸的．

记 $h(x)=f(x)-\bigl(f(\xi)+f'(\xi)(x-\xi)\bigr)$．则 $h''(x)>0$ 及 $h'(\xi)=h(\xi)=0$．于是函数 $h(x)$ 在 $(-\infty,\xi)$ 上严格递减，在 $(\xi,0)$ 上严格增加．任取 $a\in(-\infty,\xi)$，$b\in(\xi,0)$．则 $h(a)>0$，$h(b)>0$．取 $c\in\bigl(0,\min\{h(a),h(b)\}\bigr)$，则存在 $x_1\in(a,\xi)$，$x_2\in(\xi,b)$ 使得 $h(x_1)=h(x_2)=c$．

于是
\begin{align*}
  0=h(x_2)-h(x_1)&=f(x_2)-\bigl(f(\xi)+f'(\xi)(x_2-\xi)\bigr)-f(x_1)+f(\xi)+f'(\xi)(x_1-\xi)\\
  &=f(x_2)-f(x_1)-f'(\xi)(x_2-x_1).
\end{align*}
于是
\[
  f'(\xi)=\frac{f(x_2)-f(x_1)}{x_2-x_1}.
\]
\end{solution}
